\scanpage{25}
\mspage{23}
nu'', dont la variation doit refléter fidèlement les variations de structure du motif donnant naissance au cristal envisagé. Pour arriver à préciser ce dernier point, il faudrait des fondements un peu plus fermes de la théorie des motifs, comme de celle (certainement beaucoup plus élémentaire) des cristaux et de la cohomologie de De Rham. Une autre généralisation, (suggérée par comparaison du 3º avec S\unsure{e|E}rre-Tate, disant que si $S$ est de caractéristique $p$\inserted{$>0$}, alors étendre $A$ à $T$\inserted{,} c'est pareil qu'étendre le groupe p-divisible correspondant), concerne la théorie des groupes formels, dont il sera question au Chapitre 3. Pour préparer le terrain, je vais présenter d'une façon un peu différente le cristal $\operatorname{DR}^1(A/\unsure{S|A})$ associé à un schéma abélien, en utilisant explicitement la structure de groupe dudit.

2.6. De façon\deleted{\unsure{n|s}} générale, paraphrasant sur une base quelconque une vieille construction de Serre (c'est de lui que je l'ai apprise, du moins), on trouve que pour tout schéma abélien $A$ sur une base $S$, il y a une extension naturelle de $A$ par le fibré vectoriel $V(R^1f_*(\underline{O}_A))=V(\underline{t}_A)$, (où $A$ est le schéma abélien dual de $A$). Attention à la notation, le fibré vectoriel $V(\underline{E})$ est contravariant en $\underline{E}$, ces sections sont les homomorphismes $\underline{E}\longrightarrow\underline{O}_S$ ! L'extension en question est universelle parmi les extensions de $A$ par des fibrés vectoriels, est fonctorielle en $A$, et compatible avec changement de base. Appelons la $G(A)$. Ainsi, le faisceau de Lie de $G(A)$ est une extension
\diagram{\begin{tikzcd}[ampersand replacement=\&]
0 \arrow[r] \& R^1f_*(\underline{O}_{\unsure{S|A}})^{\vee} \arrow[r] \arrow[d,equal] \& \underline{g}(A) \arrow[r] \& R^0f_*(\underline{\Omega}^1_{A/S})^{\vee} \arrow[r] \arrow[d,equal] \& 0 \\
 \& \underline{t}_{A^*}^{\vee} \& \& \underline{t}_A \&
\end{tikzcd}}
,
qui est duale de la suite exacte envisagée dans 2.5., dont nous avons donc ici une construction indépendante en termes d'extensions de $A$ par des groupes vectoriels. On peut en profiter pour préciser encore les assertions 1º et 3º de 2.4. et 2.5. L'assertion 1º se précise en affirmant que, pour un prolongement $B$ de $A$, de $S$ à $T$, le schéma en groupes $G(B/S)$ est déterminé, à isomorphisme canonique \inserted{(et même \uline{unique})} près, par la seule donnée de $A$ et de \deleted{l'extension} $S\longrightarrow T$, indépendamment du choix particulier de prolongement $B$. En d'autres termes, on trouve un \uline{cristal}
\margin{24}
